% ==============================================================================
% ПАРТИЯ B03-017: DISS-005-B0735 .. DISS-005-B0754
% Глава 1: Выводы по главе 1
% Замечания: нет
% Иллюстрации: нет
% ==============================================================================

% DISS-005-B0735
\section{Выводы по главе 1}
\label{sec:ch1_conclusions}

% DISS-005-B0736
Был проведён системный анализ литературы и руководств эксплуатации различных производителей. Существующие

% DISS-005-B0737
!!Слова о том, как обстоят дела с БАВР, и т.д. Сейчас тяжко идёт, да и части пока нет, надо будет потом написать!!

% DISS-005-B0739
В результате анализа литературы и реальных осциллограмм установлено, что устройства БАВР имеют следующие проблемы функционирования и задачи, требующие дальнейшей исследований:

% DISS-005-B0740 .. DISS-005-B0750
\begin{itemize}[leftmargin=*,itemsep=0pt]
    \item Проблема затянутого срабатывания БАВР при наличии мощных электродвигателей. Суть проблемы заключается в том, что СД способны поддерживать напряжение на секциях шин на достаточном уровне длительное время, при этом, введение ускорений и модернизации пусковых органов (пуск по напряжению, пуск по углу и пуск по частоте) не всегда оказывается достаточным, так как работа БАВР блокируется поведением РНМ. В алгоритме БАВР серии МИР производства ООО «АПС» проблемы функционирования РНМ решаются путём введения адаптивных уставок по модулю тока, по углу нагрузки и по ширине зоны срабатывания. Однако, на ряде объектов требуется подстройка указанных параметров, о чём будет сказано далее.
    \item Проблема излишнего срабатывания БАВР при пуске мощных электродвигателей. Суть проблемы заключается в том, что при пуске мощных СД на текущей подстанции и ряде смежных происходит просадка напряжений, которая при некорректном настройки РНМ или отсутствии нагрузки может приводить к срабатыванию БАВР. Возможными путями решения являются:
    \begin{itemize}[leftmargin=*,itemsep=0pt]
        \item ручная подстройка адаптивных уставок (необходима высокая квалификация наладчиков);
        \item моделирование энергообъекта и выбор уставок на ней (необходима высокая квалификация проектировщика);
        \item моделирование энергообъекта, разработка алгоритма автоматической настройки БАВР на модели сети (проектировщик только создаёт модель сети, а режимы и точки КЗ выбираются автоматически, по ним автоматом рассчитываются параметры срабатывания -- подобные системы отсутствуют на данный момент);
        \item моделирование энергообъекта и машинное обучение под конкретный энергообъект (данный подход слабо изучен в электроэнергетике);
        \item машинное обучение под универсальный энергообъект.
    \end{itemize}
    \item Проблема возможности одновременного решения всех проблем путём применения жёсткой логики. Даже с учётом всех имеющихся адаптивностей, остаются проблемы, которые ввиду особенностей стечения событий не могут быть решены исключительно за счёт подстройки существующих органов РНМ в рамках жёстко программируемой логики. И, следовательно, считаем актуальным проведение научных исследований на тему применения гибкой логики, основанной на методах машинного обучения.
    \item Задача исследования поведения БАВР с новыми типами датчиков тока и напряжения. В последние годы всё чаще начинают внедряться новые датчики тока (катушки Роговского) и датчики напряжения (ёмкостные делители). При этом вопросы их применения в устройствах БАВР на данный момент остаются не исследованными и требуется проведение дополнительных исследований поведение вектора РНМ на основе указанных устройств в переходных процессах.
    \item Задача разработки органов РНМ на основе напряжений тока и напряжений нулевой последовательности. В существующих на данный момент устройствах БАВР органы регистрирующие ОЗЗ выполняют исключительно блокирующую функцию. При этом, не осуществляется локализация ОЗЗ -- внутреннее оно, или внешние. С внедрением же БАВР, предусматривающих пуск по факту регистрации однофазного замыкания, вопрос добавление органов, осуществляющих локализацию места ОЗЗ, является очень важным. И требуется дополнительное исследование данного направления.
    \item Задача широкого статистического анализа поведения РНМ. Описанные выше особенности и проблемы функционирования РНМ основаны на опыте использования и наладки устройств БАВР, а также на получении обратной связи от эксплуатации с просьбами решения возникающих сложностей. При этом, для более качественной оценки, требуется провести статистический анализ максимально широкого спектра случаев поведения, для выявления возможных проблем функционирования и полноценного сравнения используемых алгоритмов.
\end{itemize}

% DISS-005-B0752
!!НАПИСАТЬ КАКИ ЗАДАЧИ ДАЛЕЕ БУДУТ РАСМОТРЕНЫ И КАК

% DISS-005-B0753
ХОТЯ, ВРОДЕ ОНО И ВЫШЕ УКАЗАНО, НО НАДО ПРОВЕРИТЬ!!

\clearpage
